\documentclass[11pt]{article}
\usepackage{amsmath, amssymb, amsthm, mathtools, hyperref}

\newtheorem{theorem}{Theorem}
\newtheorem{lemma}[theorem]{Lemma}
\newtheorem{proposition}[theorem]{Proposition}
\theoremstyle{definition}
\newtheorem{definition}{Definition}
\theoremstyle{remark}
\newtheorem*{remark}{Remark}

\title{Counting and Combinatorics}
\author{math-foundations / concept 02}
\date{}

\begin{document}
\maketitle

\section*{Introduction}

This lesson develops the basic toolkit of finite combinatorics:
permutations, combinations, multinomial coefficients, the pigeonhole
principle, and inclusion--exclusion. The centerpiece is the binomial
theorem, which we state and prove by induction. Combinatorial reasoning
underlies all of discrete probability and many counting arguments in
machine learning theory (sample-space sizes, VC growth functions,
counting of hypothesis classes).

\section*{Basic Principles}

\begin{definition}[Multiplication Principle]
If a procedure consists of $k$ stages, with $n_i$ choices available at
stage $i$ independently of previous choices, then the total number of
ways to execute the procedure is $\prod_{i=1}^k n_i$.
\end{definition}

\begin{definition}[Addition Principle]
If a set $S$ is the disjoint union $S = \bigsqcup_{i=1}^k S_i$, then
$|S| = \sum_{i=1}^k |S_i|$.
\end{definition}

\section*{Permutations and Combinations}

\begin{definition}[Permutation]
Let $n,k$ be nonnegative integers with $k \leq n$. A $k$-\emph{permutation}
of an $n$-element set $S$ is an injection $\{1,\ldots,k\} \to S$,
equivalently an ordered tuple of $k$ distinct elements of $S$. The number
of $k$-permutations is
\[
  P(n,k) = \frac{n!}{(n-k)!} = n(n-1)\cdots(n-k+1).
\]
\end{definition}

\begin{definition}[Combination]
A $k$-\emph{combination} of an $n$-element set $S$ is a $k$-element subset
of $S$. The number of $k$-combinations is the binomial coefficient
\[
  \binom{n}{k} = \frac{n!}{k!(n-k)!}.
\]
\end{definition}

\begin{proposition}[Pascal's rule]
For $1 \leq k \leq n-1$,
\[
  \binom{n}{k} = \binom{n-1}{k-1} + \binom{n-1}{k}.
\]
\end{proposition}

\begin{proof}
Fix an element $x \in S$ with $|S|=n$. A $k$-subset either contains $x$
(there are $\binom{n-1}{k-1}$ such, choosing the remaining $k-1$ from
$S\setminus\{x\}$) or excludes $x$ (there are $\binom{n-1}{k}$). These
cases are disjoint, so the count is the sum.
\end{proof}

\section*{The Pigeonhole Principle}

\begin{definition}[Pigeonhole Principle]
If $f \colon X \to Y$ is a function between finite sets with $|X| > |Y|$,
then $f$ is not injective: there exist $x_1 \neq x_2$ in $X$ with
$f(x_1) = f(x_2)$. More generally, if $|X| > k \cdot |Y|$, then some
$y \in Y$ has at least $k+1$ preimages.
\end{definition}

\section*{The Binomial Theorem}

\begin{theorem}[Binomial Theorem]\label{thm:binomial}
For all $n \in \mathbb{Z}_{\geq 0}$ and all $x,y$ in a commutative ring,
\[
  (x+y)^n \;=\; \sum_{k=0}^{n} \binom{n}{k} x^{k} y^{\,n-k}.
\]
\end{theorem}

\begin{proof}
We proceed by induction on $n$.

\emph{Base case ($n=0$):} Both sides equal $1$. Indeed
$(x+y)^0 = 1$ and $\sum_{k=0}^{0}\binom{0}{k}x^k y^{-k}$ has its single
term $\binom{0}{0}x^0 y^0 = 1$.

\emph{Inductive step:} Assume the formula holds for some $n \geq 0$. Then
\begin{align*}
(x+y)^{n+1}
  &= (x+y)\,(x+y)^n
  = (x+y)\sum_{k=0}^{n} \binom{n}{k} x^k y^{n-k} \\
  &= \sum_{k=0}^{n} \binom{n}{k} x^{k+1} y^{n-k}
   + \sum_{k=0}^{n} \binom{n}{k} x^{k} y^{n-k+1}.
\end{align*}
Reindex the first sum with $j = k+1$:
\[
  \sum_{k=0}^{n} \binom{n}{k} x^{k+1} y^{n-k}
  = \sum_{j=1}^{n+1} \binom{n}{j-1} x^{j} y^{n+1-j}.
\]
The second sum is already indexed in the form $\sum_{j=0}^{n}
\binom{n}{j} x^{j} y^{n+1-j}$. Combining and isolating the $j=0$ and
$j=n+1$ boundary terms,
\begin{align*}
(x+y)^{n+1}
  &= \binom{n}{n} x^{n+1} y^{0}
   + \sum_{j=1}^{n} \left[\binom{n}{j-1} + \binom{n}{j}\right] x^j y^{n+1-j}
   + \binom{n}{0} x^0 y^{n+1} \\
  &= x^{n+1}
   + \sum_{j=1}^{n} \binom{n+1}{j} x^j y^{n+1-j}
   + y^{n+1}
\end{align*}
by Pascal's rule. Since $\binom{n+1}{0} = \binom{n+1}{n+1} = 1$, this is
exactly $\sum_{j=0}^{n+1} \binom{n+1}{j} x^j y^{n+1-j}$, completing the
induction.
\end{proof}

\section*{Inclusion--Exclusion}

\begin{theorem}[Inclusion--Exclusion]
For finite sets $A_1,\ldots,A_n$,
\[
  \left| \bigcup_{i=1}^{n} A_i \right|
  = \sum_{\emptyset \neq S \subseteq [n]} (-1)^{|S|+1}
     \left| \bigcap_{i \in S} A_i \right|.
\]
\end{theorem}

\section*{Worked Examples}

\paragraph{(a) Permutation.} The number of ways to arrange $3$ chosen
books in order out of $7$ available is $P(7,3) = 7\cdot 6 \cdot 5 = 210$.

\paragraph{(b) Combination.} The number of $5$-card poker hands from a
standard deck is $\binom{52}{5} = 2{,}598{,}960$.

\paragraph{(c) Pigeonhole.} Among any $367$ people, at least two share a
birthday: $|X|=367 > 366 = |Y|$.

\paragraph{(d) Inclusion--Exclusion.} Integers in $\{1,\ldots,100\}$
divisible by $2$, $3$, or $5$:
\[
  50 + 33 + 20 - 16 - 10 - 6 + 3 = 74.
\]

\paragraph{(e) Binomial theorem.} Setting $x = y = 1$ in
Theorem~\ref{thm:binomial} yields $2^n = \sum_{k=0}^n \binom{n}{k}$, the
count of all subsets of an $n$-set.

\end{document}
